% LaTeX report for the suffix-array n-gram interval-emission algorithm.
\documentclass[11pt]{article}

\usepackage[margin=1in]{geometry}
\usepackage{amsmath,amssymb,amsthm}
\usepackage{enumitem}
\usepackage{hyperref}
\usepackage{microtype}

\newtheorem{theorem}{Theorem}
\newtheorem{lemma}{Lemma}

\title{Stack-Based N-Gram Interval Emission with a Suffix Array}
\author{}
\date{}

\begin{document}

\maketitle

\section{Goal}

Let
\[
T = T[0], T[1], \ldots, T[N-1]
\]
be a tokenized corpus.  Token value \(0\) is the end-of-line sentinel
\(\mathsf{EOL}\), and a valid n-gram may not cross that sentinel.  Given a
suffix array and its adjacent longest-common-prefix array, the goal is to emit
compact records for all repeated valid n-grams whose lengths satisfy
\[
n_{\min}\leq \ell\leq n_{\max}.
\]

The report intentionally focuses on repeated n-grams.  It does not emit
singleton count-\(1\) intervals.

\section{Definitions}

\begin{description}[leftmargin=2.6cm,style=nextline]
    \item[\(T\)]
    The tokenized corpus.  The phrase ``tokens \(T[q]\) with \(a\leq q<b\)''
    means the half-open span \(T[a],T[a+1],\ldots,T[b-1]\).

    \item[\(N\)]
    The corpus length.

    \item[\(S_p\)]
    The suffix beginning at token position \(p\): the tokens \(T[q]\) with
    \(p\leq q<N\).

    \item[\(\mathit{SA}\)]
    The suffix array.  \(\mathit{SA}[i]=p\) means suffix \(S_p\) appears at
    suffix-array position \(i\).

    \item[\(\mathit{LCP}\)]
    The adjacent common-prefix array \(\mathit{LCP}[]\).  \(\mathit{LCP}[i]\) is the length of the
    longest common prefix of the suffixes at positions \(i\) and \(i+1\) in
    \(\mathit{SA}\).  Define \(\mathit{LCP}[N-1]=0\).  Entries of
    \(\mathit{LCP}[]\) stop at \(\mathsf{EOL}\), so they never describe a prefix
    crossing a line boundary.

    \item[Suffix-array interval]
    The standard suffix-array interval property is assumed: for any token
    sequence \(P\), all suffixes beginning with \(P\) occupy one contiguous
    interval in suffix-array order.
\end{description}

For a repeated interval \([L,R]\) with \(L<R\), Yamamoto and Church describe a
class of substrings by two boundary lengths:
\[
\operatorname{LBL}(L,R)=\max(\mathit{LCP}[L-1],\mathit{LCP}[R]),
\qquad
\operatorname{SIL}(L,R)=\min_{L\leq t<R}\mathit{LCP}[t],
\]
with missing boundary entries treated as \(0\).  The interval owns exactly the
lengths
\[
\operatorname{LBL}(L,R)<\ell\leq\operatorname{SIL}(L,R).
\]
Intuitively, the SIL is how far every suffix inside the interval agrees, and
the LBL is how far that agreement is already explained by a larger neighboring
interval.  The algorithm below emits each such interval-length class once.

\section{Intuition}

Read \(\mathit{LCP}[]\) from left to right as a sequence of boundaries between
adjacent suffixes.  A positive value at boundary \(b\) says that suffix-array
positions \(b\) and \(b+1\) share a positive prefix.  If later boundaries stay
high enough, that shared prefix extends across a wider suffix-array interval.
When the scan reaches a smaller value, every longer shared prefix that was
active must end there.

The stack stores those active shared prefixes.  An entry
\((\mathit{left},d)\) means: from suffix-array position \(\mathit{left}\)
through the current position \(b\), every suffix shares the same prefix of
length \(d\), and \(\mathit{left}\) is as far left as that prefix can extend.
The stored prefix lengths are strictly increasing, so shorter enclosing prefixes
are below longer nested prefixes.  When a prefix length is popped, its interval
has just ended, and the algorithm knows both its occurrence count and the range
of n-gram lengths owned by that interval.

\section{Stack-Based Interval Emission}

The algorithm emits records
\[
(\mathit{left},\mathit{right},\mathit{first\_len},\mathit{last\_len},\mathit{count}).
\]
For every length \(\ell\) with
\(\mathit{first\_len}\leq\ell\leq\mathit{last\_len}\), the n-gram represented by
tokens \(T[q]\) with
\[
\mathit{SA}[\mathit{right}]\leq q<\mathit{SA}[\mathit{right}]+\ell
\]
occurs exactly \(\mathit{count}\) times, at suffix-array positions
\(\mathit{left},\ldots,\mathit{right}\).  Thus one record may stand for several
concrete n-grams, all with the same occurrence interval and count.

\begin{figure}[p]
\scriptsize
\begin{verbatim}
stack_lcp_intervals(SA, LCP, n_min, n_max):
    stack = empty stack of (left, prefix_len)

    for each b satisfying 0 <= b < N:
        boundary_lcp = LCP[b]  # LCP[N - 1] is the final 0 sentinel.
        interval_start = b

        # Stack invariant:
        #   Prefix lengths are strictly increasing from bottom to top.
        #   Each entry (left, prefix_len) is an active interval [left, b]:
        #     every suffix at positions left, ..., b shares the same
        #     prefix of length prefix_len.
        #   The entry's left value is the earliest suffix-array position
        #     from which that common prefix is shared by every suffix
        #     through position b.
        while stack is not empty and stack.top.prefix_len > boundary_lcp:
            (left, prefix_len) = stack.pop()
            right = b

            parent_len = boundary_lcp
            if stack is not empty:
                parent_len = max(parent_len, stack.top.prefix_len)

            first_len = max(parent_len + 1, n_min)
            last_len = min(prefix_len, n_max)
            if first_len <= last_len:
                # Interval invariant:
                #   [left, right] is the maximal suffix-array interval whose
                #   suffixes share this stack entry's prefix through
                #   length prefix_len.
                #   For first_len <= ell <= last_len, the ell-token
                #   prefix has exactly this interval and no larger one.
                emit (left, right, first_len, last_len, right - left + 1)

            interval_start = left

        # A zero value starts no positive repeated-prefix interval.
        # A positive value starts, or continues, a common-prefix interval
        # beginning at interval_start and crossing boundary b.
        if boundary_lcp > 0 and
                (stack is empty or stack.top.prefix_len < boundary_lcp):
            stack.push((interval_start, boundary_lcp))
\end{verbatim}
\caption{Stack-based emission of repeated n-gram interval classes.}
\end{figure}

\section{Correctness}

\begin{lemma}[Popped intervals are maximal]
When an entry \((\mathit{left},d)\) is popped at boundary \(b\), the
suffix-array positions \([\mathit{left},b]\) form the maximal interval whose
suffixes all share that length-\(d\) prefix.
\end{lemma}

\begin{proof}
The stack invariant gives the shared-prefix part directly: every suffix at
positions \(\mathit{left},\ldots,b\) has the same prefix of length \(d\).

The entry is popped only because \(\mathit{boundary\_lcp}=\mathit{LCP}[b]<d\).
Therefore the suffix at position \(b+1\) does not share that length-\(d\) prefix, so the
interval cannot extend to the right.  The same invariant says
\(\mathit{left}\) is the earliest position from which all suffixes through
\(b\) still share that prefix, so the interval cannot extend to the left.
\end{proof}

\begin{lemma}[The emitted lengths are exactly the owned lengths]
For a popped entry with prefix length \(d\), the emitted length range
\[
\max(\mathit{parent\_len}+1,n_{\min})\leq \ell\leq
\min(d,n_{\max})
\]
contains exactly the requested valid n-gram lengths owned by the popped
interval.
\end{lemma}

\begin{proof}
All suffixes in the popped interval share prefixes up to length \(d\), so no
owned length can exceed \(d\).  Lengths at most \(\mathit{parent\_len}\) are not owned
by this interval: they are still shared by an enclosing interval, either the new
top of the stack or the interval crossing the current boundary.  Emitting them
here would duplicate that enclosing class.

Thus the unfiltered owned lengths are precisely
\(\mathit{parent\_len}<\ell\leq d\).  The lower bound \(n_{\min}\), upper bound
\(n_{\max}\), and the assumption that \(\mathit{LCP}[]\) stops at
\(\mathsf{EOL}\) then remove every invalid length.  The remaining range is
exactly what the algorithm emits.
\end{proof}

\begin{theorem}[Correctness of stack-based emission]
The algorithm emits every repeated valid n-gram class exactly once, with the
exact occurrence count for every length in the emitted range.
\end{theorem}

\begin{proof}
Let \(P\) be any valid n-gram that occurs at least twice, let
\(\ell=|P|\), and let \([L,R]\) be its maximal suffix-array interval.  By the
suffix-array interval property, suffixes \(L,\ldots,R\) all begin with \(P\),
and neither neighboring suffix outside the interval begins with \(P\).

During the left-to-right scan, the common prefix \(P\) becomes active at
boundary \(L\), remains active through boundary \(R-1\), and must end at the
first boundary after \(R\).  By the first lemma, the popped interval is
the maximal interval \([L,R]\).  By the second lemma, length \(\ell\) is emitted
for this interval exactly when it is not already owned by a child interval or an
enclosing interval.  Hence \(P\)'s class is emitted once and only once.

The emitted count is \(R-L+1\), the number of suffixes in the maximal interval.
Each such suffix is one occurrence of \(P\), so the count is exact.
\end{proof}

\section{Analysis}

In the stack pass, each active common-prefix length is pushed at most once and
popped at most once, so stack maintenance is \(O(N)\).  The total running time is
\[
O(N+Q),
\]
where \(Q\) is the number of emitted class records after applying \(n_{\min}\)
and \(n_{\max}\).  The line-boundary limit is already enforced by the definition
of \(\mathit{LCP}[]\).  If a later consumer expands each class record into
individual n-gram lengths, that expansion is output-sensitive and adds exactly
the number of concrete n-grams produced.

\section{References}

Mikio Yamamoto and Kenneth W. Church. ``Using Suffix Arrays to Compute Term
Frequency and Document Frequency for All Substrings in a Corpus.''

Kumiko Umemura and Kenneth W. Church. ``Substring Statistics.''

\end{document}
